\part{The gluing of tubes}
\section{A.1. The corrected gluing statement}
\subsection{Conventions and the distinction between three categories}
We write $\Sh(Z)$ for the Grothendieck topos of sheaves of sets on a
topological space $Z$. Geometric morphisms point in the direction of maps of
spaces; their inverse-image functors point in the opposite direction.
Throughout this part, a colimit of topoi means a bicategorical colimit,
with invertible comparison 2-cells, in $\Topos$.
For a closed subset $Y\subset Z$, the notation
\[
 \mathcal N_Z(Y)=(\Sh(N))_{N\supset Y,\ N\text{ open in }Z}
\]
denotes a \emph{pro-topos}, with transition morphisms induced by inclusions
of smaller neighborhoods. It is not an object of $\Topos$. We use
$2\Pro(\Topos)$ with mapping categories
\[
 \Hom_{2\Pro}( (E_a)_a,(T_b)_b)
 \simeq \mathop{\mathrm{bilim}}_b
       \mathop{\mathrm{bicolim}}_a\Hom_{\Topos}(E_a,T_b).
\tag{A.1.1}\label{eq:prohom}
\]
The second colimit is filtered. Constant pro-objects are denoted by $cE$.
Finally, the infinity-topos of sheaves of spaces is denoted by
$\Sh_\infty(Z)$. Its shape is discussed in Part B; it must not be identified
with a functor preserving arbitrary bicategorical colimits of ordinary topoi.

\begin{result}[A1.1; Proved: the neighborhood-limit collapse]\label{A1.1}
Suppose that $Z$ is a regular Hausdorff space and $Y\subset Z$ is closed.
The inverse limit in $\Topos$ of the open subtopoi $\Sh(N)\subset\Sh(Z)$,
for $N$ ranging over all neighborhoods of $Y$, is $\Sh(Y)$. The analogous
limit of $\Sh(N\setminus Y)$ is the empty topos. Thus taking this limit
before gluing discards precisely the punctured-neighborhood information.
\end{result}
\begin{proof}
A geometric morphism $f:E\to\Sh(Z)$ factors through $\Sh(N)$ exactly when
$f^*N=1_E$, where $N$ is regarded as a subterminal object. Regularity gives
\[
 Z\setminus Y=\bigcup_{N\supset Y}\operatorname{int}(Z\setminus N).
\]
Indeed, for $z\notin Y$ choose an open neighborhood $V$ of $z$ whose closure
avoids $Y$, and put $N=Z\setminus\overline V$. Then $N$ contains $Y$ and
$V\subset\operatorname{int}(Z\setminus N)$.
If $f$ factors through every $N$, then
$f^*\operatorname{int}(Z\setminus N)=0_E$ for every $N$. Since $f^*$
preserves unions, $f^*(Z\setminus Y)=0_E$. This is exactly the condition for
factorization through the complementary closed subtopos $\Sh(Y)$.
Factorizations and their 2-morphisms are unique in the required fully
faithful sense, so this is the limit universal property. Apply the same
argument over $Z\setminus Y$ to the opens $N\setminus Y$; their excluded
interiors cover the whole space. The resulting factorization is through
the empty subtopos. This rechecks the substantive assertion of R8, under
the stated separation hypothesis; paracompactness alone without separation
is not what this proof uses.
\end{proof}

We fix the manuscript's axioms a)--d): the closed subsets $(X_i)_{i\in I}$
are locally finite, $i\leq j$ implies $X_i\subset X_j$,
$X=\bigcup_iX_i$, and
\[
 X_i\cap X_j=\bigcup_{k\leq i,j}X_k.
\tag{A.1.2}\label{eq:abcd}
\]
Here $I$ is a poset. For base change it is essential to use the order-based
closed subset
\[
 R_i^{\mathrm{ord}}=\bigcup_{k:\,i\nleq k}X_k,
 \qquad S_i=X_i\setminus\bigcup_{k<i}X_k.
\tag{A.1.3}\label{eq:Riord}
\]
When all $S_i$ are nonempty, this agrees with the support-based $R_i$ of
151/057. When empty strata or coincident closed subsets occur, it need not
agree; we do not silently replace one definition by the other.

\begin{result}[A1.2; Proved: open-tube descent, including infinitely many strata]\label{A1.2}
For every tuple $\alpha=(N_i)_{i\in I}$ of open neighborhoods $N_i\supset X_i$
in $X$, put
\[
 W_i^\alpha=N_i\setminus R_i^{\mathrm{ord}},\qquad
 C_q^\alpha=\Sh\left(\coprod_{(i_0,\ldots,i_q)\in I^{q+1}}
              W_{i_0}^\alpha\cap\cdots\cap W_{i_q}^\alpha\right).
\tag{A.1.4}\label{eq:tubecech}
\]
The $W_i^\alpha$ form an open cover of $X$. Their augmented Cech diagram
has the following properties.
\begin{enumerate}
\item For every fixed $\alpha$, its truncation in degrees $0,1,2$ has
      $\Sh(X)$ as its bicolimit in $\Topos$.
\item Regard each $C_q=(C_q^\alpha)_\alpha$ as a \emph{single synchronized
      pro-object}. Then
      \[
        \mathop{\mathrm{bicolim}}_{[q]\in\Delta_{\leq2}^{\mathrm{op}}}C_q
        \simeq c\Sh(X)
        \quad\text{in }2\Pro(\Topos).
      \tag{A.1.5}\label{eq:procech}
      \]
\item If $I$ is finite, one can equivalently use the individual nonempty
      intersections $W_A=(\Sh(\bigcap_{i\in A}W_i^\alpha))_\alpha$ for
      $1\leq |A|\leq3$, with all their face inclusions. For infinite $I$,
      the synchronized pro-coproducts in (\ref{eq:tubecech}) are part of the
      assertion; no interchange of an infinite product and a filtered
      colimit is asserted.
\end{enumerate}
No Thom--Mather, Whitney, or homotopy-link assumption is needed. Such
assumptions are needed for geometric models and homotopy replacements of
the intersections, not for this descent theorem.
\end{result}
\begin{proof}
For $x\in X$, the finite set $I_x=\{i:x\in X_i\}$ is nonempty and downward
directed by (\ref{eq:abcd}). It has a least element $i(x)$. If $x\in X_k$,
then $i(x)\leq k$; consequently $x\notin R_{i(x)}^{\mathrm{ord}}$ and
$x\in W_{i(x)}^\alpha$. Local finiteness makes every union in
(\ref{eq:Riord}) closed, so all these sets are open.

We give the full descent argument. For an arbitrary open cover $(W_i)$,
a sheaf $F$ on $X$ restricts to sheaves $F_i$ on $W_i$ and isomorphisms
$g_{ij}:F_i|_{W_i\cap W_j}\to F_j|_{W_i\cap W_j}$, with
$g_{ii}=1$ and $g_{jk}g_{ij}=g_{ik}$ on triple intersections. Conversely,
given these data, define, for each open $O\subset X$,
\[
 F(O)=\left\{(s_i)_i\in\prod_i F_i(O\cap W_i):
             g_{ij}(s_i|)=s_j|\text{ for all }i,j\right\}.
\]
The sheaf axioms for the $F_i$ prove both existence and uniqueness of
gluing in this formula. On $W_i$, the cocycle condition identifies this
sheaf with $F_i$. The identical construction for morphisms proves an
equivalence of categories, not merely a bijection on isomorphism classes.

To obtain the topos universal property, let $T$ be any topos and suppose
geometric morphisms $f_i:\Sh(W_i)\to T$ are identified on overlaps with
the stated cocycle. Apply sheaf descent to $f_i^*(A)$ for every $A\in T$,
and to their arrows, to obtain $f^*:T\to\Sh(X)$. This functor preserves
finite limits and all colimits because restriction to the open cover does,
and because those restrictions are jointly conservative. More explicitly,
a right adjoint exists: for $F\in\Sh(X)$, take in $T$ the equalizer of
the two arrows from $\prod_i f_{i*}(F|_{W_i})$ to
$\prod_{i,j}f_{ij*}(F|_{W_i\cap W_j})$. The adjunction follows from the
displayed formula for sheaf gluing. Natural transformations between the
$f_i^*$ glue in the same way. This proves the bicolimit universal property.
Degrees $0,1,2$ encode respectively objects, overlap isomorphisms, and
their cocycle; no further degree is used.

The poset of tuples $\alpha$, ordered by componentwise shrinking, is
cofiltered, even when $I$ is infinite. The category
$\Delta_{\leq2}^{\mathrm{op}}$ has finitely many objects and morphisms.
Filtered bicolimits of categories commute with finite bilimits: every
object, isomorphism, and equation in a finite cone is represented at one
common stage, and equality of two such finite systems is witnessed after
a common refinement. This proves essential surjectivity and full
faithfulness of the comparison. Applying this fact to
(\ref{eq:prohom}), and commuting limits with limits, gives, for any
pro-topos $(T_b)_b$,
\begin{align*}
 \mathop{\mathrm{bilim}}_{q\leq2}\Hom_{2\Pro}(C_q,(T_b)_b)
 &\simeq\mathop{\mathrm{bilim}}_b
     \mathop{\mathrm{bicolim}}_\alpha
     \mathop{\mathrm{bilim}}_{q\leq2}\Hom(C_q^\alpha,T_b)\\
 &\simeq\mathop{\mathrm{bilim}}_b\Hom(\Sh(X),T_b).
\end{align*}
The last equivalences are compatible with all refinements because their
augmentations are the identity on $X$. This is
(\ref{eq:procech}). For finite $I$, descent data on the finite intersection
poset of subsets of cardinality at most three are the same data already
written out, proving the final assertion.
\end{proof}

The empty case has $I=\varnothing$ and $X=\varnothing$ and gives the
initial topos; a one-stratum cover gives the identity diagram. Empty
indexed strata are allowed throughout. Local finiteness cannot simply
be dropped: take the real line with one closed singleton $X_x=\{x\}$
for every real $x$, indexed discretely. The cover, monotonicity and
intersection axioms still hold, but the resulting $W_x=\{x\}$ are not
open. The closed-piece colimit has the discrete topology and is not
the usual real line. Thus both the openness step here and the closed
quotient step of A4.1 have a concrete failure without that hypothesis.

The theorem supplies a precise correction of 151/069, but its intersections
are additional stones. They are not, in general, any of the
$V_{ij},V_{ij}^*,V_{i,k}^j$ of 151/057--061. When $X$ is not already one of
the closed $X_i$, the global tubes used in (\ref{eq:tubecech}) amount to
adjoining an ambient top index with closed subset $X$; its open stratum may
be empty. The following counterexample proves that this extra intersection
data cannot be replaced by simply allowing all pairs and triples in the
manuscript's formulas.

\subsection{The successor and all-pair counterexamples}
\begin{result}[A1.3; Refuted: 151/069, and Lead L1's all-pair repair]\label{A1.3}
The successor-pair diagram of 151/069 does not have $\Sh(X)$ as its
colimit. Even the diagram of \emph{all} manuscript pair and triple stones,
with \emph{all} their canonical inclusions over $X$, does not suffice.
Both assertions fail for the Whitney and conically smooth stratification
\[
 X=\mathbb C,\qquad X_0=\{0\},\quad X_1=\mathbb R,\quad X_2=\mathbb C.
\]
Disconnected open strata are allowed, as in the question. The all-pair
and all-triple failure also has a connected-stratum example in
$\{0\}\subset\mathbb R^2\subset\mathbb R^4$, described below.
\end{result}
\begin{proof}
The five connected cells are $o=\{0\}$, the positive and negative rays
$p,n$, and the upper and lower half-planes $u,d$. Their incidence poset
has relations
\[
 o<p<u,\quad o<n<u,\quad o<p<d,\quad o<n<d.
\]
Every covariant set-valued functor on this poset yields a sheaf on the
plane with those cell stalks and generization maps: take its sheaf on the
finite Alexandrov poset and pull it back along the continuous cell map.
The same construction applies to the punctured plane and the real line.
This explicitly realizes all sheaves used below.

For the successor version assign a singleton to $o,p,n,d$ and the set
$\{0,1\}$ to $u$. On $V_{01}$ use singleton stalks and their unique maps.
On $V_{12}$ use the sheaf on the punctured plane whose maps $p\to u$ and
$n\to u$ have values $0$ and $1$, respectively. Its other nonidentity
maps have singleton targets. These data agree on all successor stones.
Any sheaf on the plane restricting to them would have a section
representing the unique stalk element at $o$. After shrinking to a disk,
its restriction to the connected upper half-disk must have a single value.
Restrictions near the positive and negative rays force that value to be
both $0$ and $1$. This is impossible. This independently verifies the
counterexample recorded as SA10-C3-14 and SA10-S3-K3.

For the stronger counterexample retain the same stalk sets, but put value
$1$ on \emph{both} maps $p,n\to u$ in $V_{12}$. On $V_{02}$ put value
$0$ on both maps, and also on $o\to u$. All remaining maps to singleton
stalks are unique. We list the nine distinct stones; writing
$\operatorname{germ}_Y Z$ abbreviates their pro-neighborhood systems:
\[
\begin{array}{c|c@{\qquad}c|c}
S_0&\{0\}&S_1&\mathbb R^*\\
S_2&\mathbb C\setminus\mathbb R&V_{01}&\operatorname{germ}_0\mathbb R\\
V_{01}^*&(\operatorname{germ}_0\mathbb R)\setminus\{0\}
&V_{12}&(\operatorname{germ}_{\mathbb R}\mathbb C)\setminus\{0\}\\
V_{12}^*&(\operatorname{germ}_{\mathbb R}\mathbb C)\setminus\mathbb R
&V_{02}&\operatorname{germ}_0\mathbb C\\
V_{02}^*&(\operatorname{germ}_0\mathbb C)\setminus\mathbb R&&
\end{array}
\]
Besides the standard three arrows for each pair, the only new generating
inclusions are $V_{01}\to V_{02}$ and
$V_{02}^*\to V_{12}^*$. Composites include
$V_{01}^*\to V_{12}$ and $V_{01}^*\to V_{02}$. To check completeness of
this list, a germ containing $0$ cannot map to a stone deleting $0$;
real and nonreal supports cannot contain one another; and a germ along
the whole real line cannot be contained in an arbitrarily small disk at
$0$. The remaining possibilities are exactly those listed. In particular,
$V_{12}\not\subset V_{02}$. The inclusion
$V_{02}^*\subset V_{12}^*$ does hold, because every neighborhood of the
whole real line contains a disk about $0$.

All listed restrictions agree identically on stalks and on every map they
see. For the unique strict triple $0<1<2$, formula (11) of 151/060 gives
$S_{0,2}=S_{1,2}=\mathbb R$, hence
$V_{0,2}^{0}=V_{0,2}^{1}=V_{02}^*$ and $V_{0,2}^{2}=V_{02}$.
Triples with repeated indices reproduce pair stones. They impose no new
condition. But the actual overlap of $V_{02}$ and $V_{12}$ contains the
punctured disk germ, where the two $p\to u$ maps differ. If a global sheaf
existed, its identifications on $S_1,S_2$ and on these overlaps would make
the two maps identical. The discrepancy survives every shrinking, so no
global sheaf exists.

We checked the primary scan 151/060 as well as its transcription. In
formula (10), the overwritten first index of $R$ is potentially read as
$i$ or $j$; the following sentence specifies $R_{i,k}$. Neither reading
repairs this counterexample. For the only interior index $j=1$, the
deletion of $S_{1,2}=\mathbb R$ absorbs both $R_{0,2}=\varnothing$ and
$R_{1,2}=\{0\}$. At the upper endpoint, reading $R_{j,k}$ instead
merely replaces the repeated $V_{02}$ by $V_{02}^*$, both already in
the pair diagram. The doubtful letter $S$ versus $\Sigma$ denotes the
set explicitly defined in (11). The later doubtful bound in the
pseudo-partition is not used: (11) and its adjacent $j=k$ note specify
the strict upper bound. If that bound were changed to $l\leq k$, its
extra $X_k$ would make the corresponding triple stone empty; this
weaker diagram still cannot detect the missing overlap.

This sheaf obstruction tests a genuine topos colimit universal property.
The object-classifying topos $\mathbf{Set}^{\mathrm{FinSet}}$ has
$\Hom(E,\mathbf{Set}^{\mathrm{FinSet}})\simeq E$: a left-exact functor
$\mathrm{FinSet}^{\mathrm{op}}\to E$ is $K\mapsto A^K$ for a unique
object $A$, and the presheaf-topos universal property gives the claimed
equivalence. Thus the compatible sheaves above are a compatible cocone of
geometric morphisms to this classifier, with no extension from $\Sh(X)$.
The argument applies unchanged in the pro-bicategory, using representatives
and their germs. For a version with every open stratum connected, replace
the chain by $\{0\}\subset\mathbb R^2\subset\mathbb R^4$. On the closed
plane use the singleton sheaf, and on its complement use the constant
two-point sheaf. The two constant choices $0,1$ define two global
closed--open recollement sheaves $F^0,F^1$ by B1.1. Put $F^0$ on $V_{02}$
and $F^1$ on $V_{12}$. Their restrictions to all named strata and to
$V_{01}$ agree, whereas their maps from the middle stratum to the top
one disagree on every punctured-origin overlap. The list of possible
germ inclusions and the equality $V_{02}^{0}=V_{02}^{1}=V_{02}^*$ are
unchanged. The strata are connected smooth manifolds; the nested linear
subspaces satisfy Whitney's conditions and have their product conical
control data. Thus connectedness does not repair the all-pair claim.
The program \texttt{check\_A\_finite.py} verifies the
finite functor equations and all 14 nonidentity inclusions in the
all-pair example. The analytic realization just given verifies the
topological hypotheses, which the finite program alone could not verify.
\end{proof}

\begin{result}[A1.4; Proved: precise universal minimality]\label{A1.4}
For ordinary-topos open-cover descent, vertices, pair overlaps, and triple
cocycles are the smallest universal Cech arities. In the subset-indexed
presentation of A1.2(3), each vertex, pair, or triple intersection object
can be indispensable for an appropriate finite open cover;
intersections of four or more members are unnecessary. There is no
claim that every slot of every particular cover is indispensable.
Ordered repetitions and degeneracies in the full Cech diagram are not
being counted as independently necessary pieces.
In binary gluing, the three-piece span
$\Sh(N)\leftarrow\Sh(N\cap U)\to\Sh(U)$ is universally indispensable;
an extra soul $\Sh(Y)\to\Sh(N)$ with no other incidence is redundant.
\end{result}
\begin{proof}
The sufficiency is the explicit descent proof of A1.2. To see necessity
of a vertex slot, let a point belong only to that member of a cover;
deleting it removes all information at the point. For a pair slot, let
only those two cover members contain a point. Without their intersection,
one may assign stalk sets of cardinalities one and two at the two copies,
which cannot come from a single sheaf. For a triple slot, let exactly the
three members contain a point and assign a two-element set to each.
Take two pair identifications to be the identity and the third to be the
transposition. Every inverse rule holds, but the triple cocycle fails.
Other cover members may be taken disjoint from that point. These are
finite counterexamples; the transposition calculation is in the program.
The binary examples are the same missing-overlap example and the missing
vertex examples, realized, if desired, by a tube and its complement in
the real line. An object with its sole outgoing arrow into an already
present object contributes no additional bicolimit data: its cocone map
is forced by composition, with a contractible choice of comparison.
That proves the assertion about the soul and also explains why absolute
minimality of the manuscript's named list cannot be demanded.
\end{proof}

\subsection{Comparison with unzipping and SGA 4}
\begin{result}[A1.5; Proved: exact criterion for deleting open pieces]\label{A1.5}
Let $X$ be a $T_1$ space, $C$ a small category, and
$D:C\to\operatorname{Open}(X)$ a diagram of open subsets and inclusions.
For $x\in X$, let $C_x$ be the full subcategory on the objects $c$ with
$x\in D(c)$. The canonical cocone has $\Sh(X)$ as its bicolimit if and
only if the groupoid obtained by inverting every arrow of $C_x$ is
equivalent to the terminal groupoid, for every $x$. Equivalently, each
$|N(C_x)|$ is nonempty, connected, and simply connected.
This gives an exact test for any proposed deletion of objects or arrows
from a specified honest-open gluing diagram. An inclusion-minimal
subdiagram is precisely one satisfying this test for which every further
permitted deletion fails the test at some point.
\end{result}
\begin{proof}
A sheaf descent datum on $D$ is a family of sheaves $F_c$ with
isomorphisms $F_d|_{D(c)}\simeq F_c$ for every arrow $c\to d$, satisfying
the functorial identities. At $x$ its stalks therefore define a functor
from the groupoidification of $C_x$ to sets. If that groupoid is terminal,
all these stalks identify uniquely, with no possible monodromy.
Take the disjoint union of the etale spaces of the $F_c$, and quotient
by the identifications on the open subspaces specified by arrows. The
quotient map is open: the saturation of an open set is the union of its
images under finite strings of these open-subspace homeomorphisms.
At a point $x\in D(c)$, connectedness gives surjectivity from the
$c$-stalk to the quotient fibre, and trivial groupoid monodromy gives
injectivity. Thus each etale space maps homeomorphically onto the
quotient over $D(c)$. The quotient is an etale space over $X$, since
the $D(c)$ cover $X$. It gives a sheaf restricting to the specified
datum. Morphisms descend by the same quotient construction, proving an
equivalence of categories. Applying the argument to inverse-image
functors and their natural transformations, as in A1.2, proves the
topos bicolimit property.

Conversely, if $C_x$ is empty, two distinct skyscraper sheaves at the
closed point $x$ restrict identically to the entire diagram. If $C_x$
is disconnected, choose stalk sets of different cardinalities on two
components. If it is connected with nontrivial groupoid isotropy, choose
the regular action of one of its nontrivial automorphism groups and
transport it to all its objects. These choices are nonconstant descent
data in sets on $C_x$. They extend to sheaf descent data on $D$: for
$x\in D(c)$ use the skyscraper with the specified stalk, and for
$x\notin D(c)$ use the terminal sheaf. Restrictions to opens missing
$x$ are terminal, so all required identifications and equations hold.
A sheaf on $X$ would induce the constant stalk datum on $C_x$, contrary
to the choice. The object classifier in A1.3 converts this obstruction
to failure of the topos universal property. Finally, the nerve of a
category has fundamental groupoid equal to its groupoidification:
edges give arrows and formal inverses, and its triangles impose exactly
the composition relations. Higher simplices do not change that
presentation. This proves the stated equivalent topological criterion.
\end{proof}

The geometric unzipping theorem of Ayala--Francis--Tanaka
\cite[Theorem 7.3.8, Corollary 7.3.7 and Proposition 8.2.3]{AFT}
provides, under conical smoothness and their finite-dimensional hypotheses,
a resolution with corners, a compact link bundle along a deepest stratum,
and its cone neighborhood. Taking that neighborhood and its open complement
produces exactly a binary cover to which A1.2 applies. Repeating the
geometric construction retains the boundary incidences and their
overlaps; it does not identify them with the manuscript's all-pair list.
The plane example therefore refutes the asserted equivalence between
these two procedures in Lead L1.

SGA 4, Expose IV, Theorem 9.5.4 \cite{SGA4}, reconstructs a topos from an
open subtopos, its closed complement, and the functor $i^*j_*$. Its objects
are triples $(A,B,A\to i^*j_*B)$. The arrow is not required to be invertible.
Our open-cover descent uses invertible overlap identifications; the
noninvertible generization information is already present inside each
mixed tube. A1.3 shows that omitting an actual overlap fails to compare
exactly this information. The two formulations are consistent once the
neighborhood system, its overlap, and the gluing functor are retained.

\section{A.2. Base change of tubes and locality}
\begin{result}[A2.1; Proved: exact neighborhood criterion]\label{A2.1}
Let $f:Z'\to Z$ be continuous and $Y\subset Z$ closed. The canonical
comparison from the neighborhood pro-object of $f^{-1}Y$ in $Z'$ to the
pullback of $\mathcal N_Z(Y)$ is an isomorphism over $Z'$ if the family
$\{f^{-1}N:N\supset Y\text{ open}\}$ is cofinal among neighborhoods of
$f^{-1}Y$. Equivalently, for every closed $F'\subset Z'$ disjoint from
$f^{-1}Y$, there is an open $N\supset Y$ with $f^{-1}N\cap F'=\varnothing$.
This holds for every closed continuous map $f$, in particular for proper
maps between locally compact Hausdorff spaces. No local-immersion
hypothesis is necessary.
\end{result}
\begin{proof}
The two neighborhood systems have the canonical inclusion morphisms over
$Z'$. Cofinality is precisely the assertion that each neighborhood in the
first system contains a pullback neighborhood, giving inverse morphisms
of these pro-subspaces. Replacing an open neighborhood $O'$ by its closed
complement gives the equivalent condition. If $f$ is closed, then
$f(F')$ is closed and disjoint from $Y$, so $N=Z\setminus f(F')$ works.
All comparisons are induced by actual inclusions, hence the same
isomorphisms hold for their pro-topoi and after deleting compatible
closed subsets. Conversely, an inverse in the category of pro-subspaces
over $Z'$ must, for each $O'$, factor some pullback neighborhood into
$O'$; this is cofinality. This is an exact criterion for this specified
over-space comparison, not a claim about abstract isomorphisms of
unlabelled pro-objects.
\end{proof}

\begin{result}[A2.2; Refuted: unrestricted open base change and support-based $R$]\label{A2.2}
Open immersions do not generally satisfy A2.1. Moreover, formula (4) of
151/071 for the support-based $R_{ij}$ fails under arbitrary inverse
images, even for a closed immersion. The order-based definition
(\ref{eq:Riord}) does commute with every inverse image.
\end{result}
\begin{proof}
Take $Z=\mathbb R$, $Y=\{0\}$ and $Z'=\mathbb R\setminus\{0\}$.
The induced center is empty, so its neighborhood pro-object has the empty
set as a cofinal stage. The pulled-back neighborhoods are punctured
interval germs, and none is empty. They cannot factor into the empty
stage over $Z'$, so the comparison is not an isomorphism. Their pro-homotopy
types even have two components instead of zero.

For the second assertion use $X_0=\{0\}\subset X_1=\mathbb R$ and the
closed immersion $f:\{0\}\hookrightarrow\mathbb R$. Originally $R_1=X_0$.
After pullback $X'_0=X'_1=\{0\}$, so no $X'_k$ fails to contain $X'_1$;
the recomputed support-based $R'_1$ is empty. But $f^{-1}R_1=\{0\}$.
The discrepancy is removed if the index condition $i\nleq k$ is held
fixed: inverse image commutes with the union in (\ref{eq:Riord}), and
likewise with the fixed-index unions defining $S_{j,k}$.

For completeness, the order-based identity is always
$R_i^{\mathrm{ord}}\cap X_i=\bigcup_{h<i}X_h$. Indeed, intersect each
$X_k$ with $X_i$ and use (d); a lower index $h\leq i,k$ cannot equal
$i$ when $i\nleq k$. Conversely every $h<i$ occurs in the defining
union. If $S_i$ is nonempty, the argument of A4.4 shows
$X_i\subset X_k$ exactly when $i\leq k$, recovering the manuscript's
support-based identity (4) under this hypothesis. Without it the
coincident-piece example above refutes (4). The separate parenthesis
after (12) of 151/060 is false even with nonempty strata: in the
three-point chain $X_0=\{0\}\subset X_1=\{0,1\}\subset
X_2=\{0,1,2\}$, take $(i,j,k)=(0,1,2)$. Then
$R_{1,2}\cap S_{1,2}=\{0\}$, whereas $\dot X_0=\varnothing$.
The finite identities are checked in the attached program. This
rechecks the substantive R3 qualification and its separately false
parenthesis without relying on the earlier check labels.
\end{proof}

The supplied second-pass transcription of 151/071 explicitly inserts
\emph{mais pas} before the reference to an open immersion. Thus the prompt
and \texttt{study/STUDY.md}, line 221, misstate this sentence. The
mathematical negation above is independent of that reading: it proves the
claim false if the parenthesis were read positively, and confirms the
negative reading if it is retained.

\begin{result}[A2.3; Proved: the corrected locality statement]\label{A2.3}
Descent for a specified open-cover diagram is local on $X$: restrictions
of that cover to an open cover of $X$ verify the same gluing equivalence,
and these local equivalences glue. The argument proposed and retracted
in 151/072 is invalid when it replaces the pullback of a closed-center
pro-neighborhood by the independently recomputed neighborhood system
on an open subset.
\end{result}
\begin{proof}
For honest opens, $(W_i\cap W_j)\cap O=(W_i\cap O)\cap(W_j\cap O)$,
with identical equalities for triples. The explicit sheaf construction
in A1.2 commutes with restricting its argument $O$ to any smaller open.
Therefore it is a local inverse to restriction, and inverses agree on
overlaps by uniqueness of sheaf gluing. Applying this level by level
proves locality for a specified pulled-back synchronized system.
A2.2 shows that recomputation can change even an empty center into a
nonempty punctured pro-object. Thus the missing comparison in 151/072
cannot be inferred from open base change. This refutes the inference,
while the valid local descent theorem remains as strong as stated.
\end{proof}

\section{A.3. Differences of cribles}
\begin{result}[A3.1; Proved and Refuted: exact closed restriction, false open homotopy comparison]\label{A3.1}
Let $I''\subset I'\subset I$ be lower ideals, and put
$Z=X_{I''}$, $X'=X_{I'}$, $U=X'\setminus Z$. The induced strata of $U$
are exactly the $S_i$ with $i\in I'\setminus I''$. For retained $i,j$,
passing from $X$ to the closed union $X'$ preserves the pair stones
defined inside $X_j$ (with fixed order-based deletions, or nonempty strata).
Passing to $U$ need not preserve even their number of connected
components or their homotopy types, including in the conically smooth
Whitney example of A1.3.
\end{result}
\begin{proof}
If $x\in S_i$ and $x\in X_k$ for $k\in I''$, then $i\leq k$ forces
$i\in I''$. This proves the assertion about remaining strata and the
same argument gives the induced lower unions. Since $I'$ contains every
index below each of its elements, the ambient $X_j$ and all its relevant
closed deletions are unchanged on restriction to $X'$. For support-based
deletions under nonempty strata, $X_i\subset X_k$ is equivalent to
$i\leq k$, so this is the same argument.

Now take $I'=\{0,1,2\}$ and $I''=\{0\}$ in A1.3. Originally
\[
 V_{12}=(\operatorname{germ}_{\mathbb R}\mathbb C)\setminus\{0\}.
\]
A cofinal family of neighborhoods of $\mathbb R$ consists of strips
$\{|y|<\varepsilon(x)\}$ with $\varepsilon:\mathbb R\to(0,\infty)$
continuous. Such strips are homeomorphic to $\mathbb R\times(-1,1)$,
and deleting the origin gives homotopy type $S^1$. The inclusions between
shrunk strips preserve the once-around-origin generator and are homotopy
equivalences, so the pro-homotopy type is constant $S^1$.
In $U=\mathbb C^*$ the induced center is $\mathbb R^*$.
There is a cofinal family of two disjoint variable-width strips over the
positive and negative rays, each of width less than $|x|/2$. Each
component contracts onto its ray. The induced tube therefore has the
constant homotopy type of two points. Cofinality of both strip families
follows by choosing a positive continuous minorant of the distance to
the complement, using paracompactness of the real line or its open rays.

Deleting the real center from the original strip gives two contractible
components, upper and lower. Deleting the center from the induced pair
of strips gives four contractible components. Hence the maps on deleted
tubes induce the nonbijective component map $\{1,2,3,4\}\to\{1,2\}$.
Both comparisons of 151/075 fail to be homotopy equivalences. This is a
cofinal-system calculation, not an artifact of a badly chosen representative.
\end{proof}

\begin{result}[A3.2; Proved: comparison of constructible sheaves]\label{A3.2}
For an open union of strata $j:U\hookrightarrow X'$, restriction
$j^*: \Sh_{\mathrm{c}}(X')\to\Sh_{\mathrm{c}}(U)$ is essentially
surjective and has the fully faithful extension-by-empty left adjoint
$j_!$. It is an equivalence exactly when $Z=X'\setminus U$ is empty.
Its full reconstruction is the recollement with $\Sh_{\mathrm{c}}(Z)$
and the nearby-section functor of B.1, with its constructibility conditions.
\end{result}
\begin{proof}
For sheaves of sets, $j_!$ exists for every open immersion. Its stalks
are the original stalks on $U$ and empty on $Z$, so it is constructible
for this stratification. The identity $j^*j_!=1$ proves essential
surjectivity and full faithfulness of $j_!$. If $Z\ne\varnothing$,
let $i:Z\hookrightarrow X'$ and compare the two sheaves
$i_*\underline{\{0\}}$ and $i_*\underline{\{0,1\}}$. They are
constructible: on $Z$ they are locally constant, and on $U$ both have
singleton stalks. They are not isomorphic at a point of $Z$, but their
restrictions to $U$ are isomorphic. Thus restriction cannot be an
equivalence. If $Z=\varnothing$ it is the identity. The recollement
description follows from the explicit inverse in B.1 or SGA 4 IV 9.5.4.
\end{proof}

\section{A.4. Closed diagrams, infima, and stable graphs}
\begin{result}[A4.1; Proved: the reconstruction already valid in R2]\label{A4.1}
Under a)--d), the strata $S_i$ partition $X$, and the colimit of the
closed-subset diagram $i\mapsto X_i$, equivalently the coequalizer of
its flag diagram, is $X$ as a topological space. Empty strata may be
discarded without destroying a)--d).
\end{result}
\begin{proof}
The least element $i(x)$ constructed in A1.2 is precisely the unique
index with $x\in S_{i(x)}$. Any two representatives of $x$ in $X_i$ and
$X_j$ are identified through $X_k$ for some $k\leq i,j$ containing $x$,
by (d). Thus the underlying set coequalizer is $X$. The map
$\coprod_i X_i\to X$ is closed: a closed subset of the coproduct is a
family of closed subsets of the closed $X_i$, and their union is closed
by local finiteness. It is surjective, hence a quotient map, proving the
topological assertion. Finally,
$X_k=\bigcup_{i\leq k,\,S_i\ne\varnothing}X_i$, because $i(x)\leq k$
for every $x\in X_k$. This identity verifies the cover and intersection
axioms after discarding empty strata; the other axioms are inherited.
This rechecks SA10-C3-01, -02, and -03.
The primary scan 151/070 really writes a strict lower bound in its
repeated intersection axiom. Taken literally, this forces $X$ empty
under local finiteness: for any supposed point choose its least
containing index $i(x)$ and apply that axiom to $i=j=i(x)$. Its right
side has no member containing the point. Thus all nonempty applications
here explicitly use the nonstrict axiom of 151/052, not the slip on
151/070. This verifies the mathematical correction R6.
\end{proof}

\begin{result}[A4.2; Refuted: the corollary of 151/052--053]\label{A4.2}
A finite meet-semilattice diagram of closed embeddings need not embed
its members in its colimit, even if its associated ordered space satisfies
(23a)--(23b) of 151/053.
\end{result}
\begin{proof}
Let $I$ be the seven proper subsets of $\{1,2,3\}$, including the empty
set, ordered by inclusion. Its meet is intersection. Assign the singleton
$\{*\}$ to $\varnothing$ and the discrete three-point space
$\{*,0,1\}$ to each of the other six indices. Bottom arrows send $*$ to
$*$. Every proper transition map between the other indices is the identity except
$X_{\{1\}}\to X_{\{1,2\}}$, which swaps $0,1$ and fixes $*$. Every
composable pair of strict arrows starts at the bottom and fixes $*$, so this is a
functor, and all maps are closed embeddings. The disjoint union is a
finite discrete ordered space: comparability is given by the graphs of
the transition maps. Reflexivity, antisymmetry, and transitivity hold,
and its order relation is closed. Thus (23a)--(23b) hold.
The six non-bottom indices form a hexagon. Going once around it transports
$0$ to $1$, since exactly one transition is a transposition. Hence the
nineteen points of the disjoint union have two equivalence classes in the
colimit, one for $*$ and one for $0,1$. No three-point member embeds.
Deleting the $*$ orbit gives the variant with an empty bottom, twelve
points and a singleton colimit. Both complete finite checks are in
\texttt{check\_A\_finite.py}; the first version avoids any convention
about whether an indexed partition may contain an empty part.
\end{proof}

\begin{result}[A4.3; Proved: exact and geometric repairs of the corollary]\label{A4.3}
For a finite diagram of closed embeddings indexed by a finite poset,
all maps to its topological colimit are closed embeddings if and only if
no zigzag of transition identifications identifies two distinct points
in one member. If $I$ is a finite meet-semilattice, the stronger condition
\[
 X_{i\wedge j}\xrightarrow{\sim}X_i\times_{X_k}X_j
 \quad(i,j\leq k)
\tag{A.4.1}\label{eq:cartesianmeet}
\]
is sufficient, and then the intersection of the images of $X_i,X_j$ in
the colimit is exactly the image of $X_{i\wedge j}$.
\end{result}
\begin{proof}
The zigzag condition is exactly injectivity of each set map to the
colimit. Assuming it, any zigzag joining two members can have a repeated
index removed, because a revisit to an index must revisit the same
point. There are only finitely many simple index paths. Transport of a
closed subset along each step is either a closed image under a closed
embedding or its inverse image. Thus its saturation meets each member
in a finite union of closed subsets. The quotient topology shows that
each member's map is closed, and the same argument applies to every
closed subset of it. An injective continuous closed map is an embedding.
The converse follows immediately from injectivity.

Under (\ref{eq:cartesianmeet}), follow a zigzag from a point of $X_i$.
Inductively maintain a common lower index and a point mapping both to
the starting point and to the current representative. An upward step
preserves this datum. For a downward step to $X_j$, use the pullback
square with the old lower index and $j$ inside their common upper
index, replacing the lower index by their meet. The point lifts
uniquely. At the end, representatives in $X_i$ and $X_j$ therefore come
from a lower index, and hence from $i\wedge j$. Taking $i=j$ proves the
zigzag condition. Taking arbitrary $i,j$ gives the claimed intersection
identity, including its subspace topology since all maps are closed
embeddings. This proves the repair without assuming that a poset
meet alone forces a pullback square.
\end{proof}

\begin{result}[A4.4; Proved: the class defined by (24)]\label{A4.4}
Under a)--d) with nonempty strata, the strengthened condition (24) of
151/054 is equivalent to the family of closed subsets $\{X_i\}$ being
closed under binary intersection, with the intersection represented by
one index. The conditional version, requiring infima only for pairs
with a common lower bound, is equivalent to this assertion for nonempty
intersections. The literal all-pairs version additionally needs a
bottom index representing the empty intersection whenever disjoint
closed strata occur.
\end{result}
\begin{proof}
If $X_i\subset X_j$ and $S_i\ne\varnothing$, choose $x\in S_i$.
Then $i=i(x)\leq j$. Thus $i\mapsto X_i$ is an order embedding.
An index representing $X_i\cap X_j$ is consequently precisely the
greatest lower bound. Conversely, if the infimum exists, (d) says that
the union of all lower closed subsets is its closed subset, giving the
intersection equality. Since every retained $X_k$ contains a nonempty
stratum, an intersection is nonempty exactly when a common lower
index exists. A formal empty bottom handles the remaining pairs.
Every finite meet-semilattice is realized in this general topological
class by its upper-Alexandrov space, with $X_i=\downarrow i$; so the
criterion describes exactly a closure-system condition and imposes
no hidden smoothness assertion.
\end{proof}

\begin{result}[A4.5; Proved: stable-graph classification and L4]\label{A4.5}
Let $2g-2+n>0$, and order isomorphism classes of connected stable graphs
of genus $g$ with $n$ individually labelled legs by reverse contraction:
$H\leq G$ if $H$ contracts to $G$. Every pair having a common lower
bound has an infimum exactly for
\[
 g=0\ (n\geq3),\qquad g=1\ (1\leq n\leq3),\qquad (g,n)=(2,0).
\tag{A.4.2}\label{eq:smallgn}
\]
Consequently the conditional interpretation of Lead L4 and the quoted
EXIT\_TOWER A.2 is correct. Literal (24) requires the empty-bottom
qualification of A4.4. This statement concerns graph types, not fixed
contractions or the incidence category.
\end{result}
\begin{proof}
A stable graph has vertex genera $g_v\geq0$, labelled legs, loops and
multiple edges, with $2g_v-2+\operatorname{val}(v)>0$ and
$g=\sum_vg_v+|E|-|V|+1$. For genus zero it is a tree, and its edges
determine stable bipartitions of the labelled legs. Distinct edges give
distinct bipartitions. A common refinement supplies a tree containing
the two sets of bipartitions; contract precisely the edges outside their
union. The resulting tree is independent of the common refinement:
root at one labelled leg, orient each split by the side avoiding that
leg, and reconstruct vertices from the inclusion hierarchy of these
nested or disjoint subsets. This reconstruction is unique. Every other
common refinement contracts to it, so it is the infimum.

For the remaining positive cases we give a complete finite certificate,
with a proof of exhaustive generation. The program
\texttt{stable\_graph\_meets.py} starts with the corolla, then applies
every inverse single-edge contraction: a loop lowering a positive
vertex genus by one, or a split of a vertex into two joined by an edge,
with all partitions of its genus and incident half-edges and legs.
It rejects precisely the unstable splits and identifies graphs by all
vertex permutations preserving the data. Contracting any edge of a
stable graph preserves stability; iterating contracts every graph to
the corolla. Thus induction on edge number proves that the enumeration
is exhaustive. The bound $|E|\leq3g-3+n$ follows by summing the vertex
stability inequalities, so the process terminates. The program computes
all contraction closures, all pairs, and all minimal common refinements.
For $(g,n)=(1,1),(1,2),(1,3),(2,0)$ the graph counts are respectively
$2,5,23,7$; all $1,10,253,21$ unordered pairs were checked, and none
has two minimal common refinements. The graph lists and contraction
closures are attached as finite certificates.

Here are uniform counterexamples in every other case; they also explain
the failures independently of enumeration. If $g\geq3$, let $G$ be the
one-edge separating graph with vertex genera $1,g-1$, and let $H$ be
the one-loop graph with one vertex of genus $g-1$. Put all markings,
if any, on the second vertex of $G$. One two-edge common refinement
has a loop on a genus-zero vertex joined to a genus-$(g-1)$ vertex;
the other has vertex genera $1,g-2$, with the loop on the second
vertex. Put all markings on the second vertex in both. Loop
contraction gives $G$, and bridge contraction gives $H$. Both are
stable and are nonisomorphic, since one has a genus-zero vertex and
the other does not. Their only one-edge contractions are $G,H$, which
are different. They are therefore distinct minimal common refinements.

If $g=2,n\geq1$, take $G$ with two genus-one vertices joined by an
edge, all legs on one vertex, and $H$ with one genus-one vertex,
all legs, and one loop. The refinements consist of a genus-zero loop
vertex joined to a genus-one vertex; in one refinement all legs lie
on the loop vertex, and in the other they lie on the other vertex.
They are stable, nonisomorphic, and minimal by the same two-edge
argument.

Finally let $g=1,n\geq4$. Let $G$ be a two-edge cycle with leg
partition $\{1\}\mid\{2,\ldots,n\}$, and $H$ the analogous cycle
with partition $\{2\}\mid\{1,3,\ldots,n\}$. A common refinement is
the triangle with vertex leg sets $\{1\},\{2\},\{3,\ldots,n\}$.
Another is the square whose opposite vertices carry $\{1\},\{2\}$,
and whose other vertices carry $\{3\},\{4,\ldots,n\}$. Every
vertex is rational and stable. Contracting suitable edges gives both
$G$ and $H$. In any common refinement obtained from the square,
the vertices carrying legs $1$ and $2$ must each remain separate
from all other marked vertices. Every square edge touches one of
these vertices, so no edge can be contracted while retaining both
requirements. The square is minimal. The triangle is also minimal,
since a two-edge common refinement would have to equal both $G$ and
$H$. These two nonisomorphic minimal refinements rule out an infimum.
This covers all stable $(g,n)$ not in (\ref{eq:smallgn}).

For the application to moduli spaces, the local smoothing parameters
of a stable curve correspond to its nodes, and smoothing a selected
set contracts the corresponding edges. The stable-graph stratum
closure relation is therefore exactly this order; see
\cite[Proposition 4.14, Section 6.1, and Section 7.3]{Schmitt}.
Each common refined type occurs, by gluing smooth pointed components
of its prescribed genera. Thus the combinatorial failures are actual
failures of the closure-intersection condition in the stated moduli
stratification, not merely formal graphs.
\end{proof}
